% Auxiliary Gaussian grouping. No preamble. Unique prefix ag-.
\section{Auxiliary grouping recovers a fixed quantum support weight}
\label{ag-section}

The physical averaging ratio and the size of a Gaussian expansion
cell need not coincide. We keep the exact source ratio $L$ fixed
and group its Gaussian cells by a second integer $J$. This gives
a finite-step recovery of one fixed connected-carrier exponent,
without assigning arbitrary reserve to inherited quantum activities.
It does not give contraction of activity amplitudes, stable analytic
radii, or a quantum RG induction. All constants below are uniform
in old refinement $M$ and compatible periodic volume after the
finite parameters $L,J$ have been chosen.

\subsection{The four graphs and parameter order}
\label{ag-graphs-section}

Fix the physical odd source ratio $L\ge5$. On the current unit
lattice, source blocks have side $L$ and index graph
$\Lambda_{\mathrm{src}}$. This is also the next-current unit graph
after rescaling lengths by $L$. Choose an odd auxiliary integer
$J$, conveniently $J=L^j$. Gaussian integration cells group
$J^d$ source blocks, so their current physical side is $JL$;
write their graph as $\Lambda_G$. After rescaling, these are
$J$-blocks of the next current lattice.

At every stage, inherited quantum carriers are $J$-blocks of
that stage's current unit lattice; call their graph $\Lambda_Q$.
The two projections are
\begin{equation}\label{ag-projections}
 \rho_J:\Lambda_{\mathrm{src}}\longrightarrow\Lambda_G,
 \qquad
 \pi_L:\Lambda_Q\longrightarrow\Lambda_G.
\end{equation}
Their fiber sizes are $J^d$ and $L^d$, respectively. Thus the
coarsening ratio of inherited carrier graphs remains $L$.

A centered odd $JL$ cube is both a union of $J^d$ centered
$L$ cubes and a union of $L^d$ centered $J$ cubes. The two
smaller tilings need not be nested in one another; their common
$JL$ tiling is enough. Work on a cofinal family of compatible
tori with periods divisible by $JL$. Taking $J=L^j$ preserves
the original hierarchy on sufficiently large members. No claim
for an incompatible short torus is needed.

The relevant dimensions are
\begin{equation}\label{ag-dimensions}
 D_G=J^d d(L^d-1)\dim\mathfrak g,
 \quad D_{\mathrm{bg}}=d(JL)^d\dim\mathfrak g,
 \quad D_{\mathrm{in}}=dJ^d\dim\mathfrak g.
\end{equation}
They count retained Gaussian coordinates, all background
coordinates of a Gaussian macrocell, and background coordinates
of an inherited carrier cell. Grouping is a coordinate
rearrangement. It changes no source average, pivot, or cutoff.

Choose $L$ first using the fixed weight test below. At that
fixed $L$, the previous theorems supply positive covariance,
classical precision and stationary-background decay exponents.
Then choose $J$ large enough for their finitely many conversion
inequalities. Finally choose the product radius and coupling
budgets. No large-$L$ lower bound for the old decay exponents
is required in this order.

\subsection{Covariance and precision on macro Gaussian cells}

Let the unchanged real baseline satisfy
$\lambda I\le C_0\le\Lambda_0I$ and, on source blocks,
\begin{equation}\label{ag-source-covariance}
 K_{\mathrm{src}}(\gamma_0)
 =\sup_i\sum_j e^{\gamma_0d_{\mathrm{src}}(i,j)}
                   \|(C_0)_{ij}\|<\infty,
 \qquad \gamma_0>0,
\end{equation}
with the same column bound. These constants are uniform in
$M$ and volume for the fixed source ratio $L$.
For $i$ in macrocell $A$ and $j$ in macrocell $B$,
\begin{equation}\label{ag-distance-projection}
 d_G(A,B)\le J^{-1}d_{\mathrm{src}}(i,j)+d.
\end{equation}
It follows coordinatewise from block projection, including
periodic distances. Since
$\|(C_0)_{AB}\|\le\sum_{i\in A,j\in B}\|(C_0)_{ij}\|$,
\begin{equation}\label{ag-covariance-bound}
 K_G(\gamma):=\sup_A\sum_B e^{\gamma d_G(A,B)}\|(C_0)_{AB}\|
 \le J^d e^{\gamma d}K_{\mathrm{src}}(\gamma/J).
\end{equation}
The same bound holds for columns. The spectrum is unchanged:
the macro fibers are direct sums of the original Euclidean
fibers, not averaged variables.

For any terminal set $X$, projection of a terminal spanning
tree gives
\begin{align}
 \tau_G(\rho_JX)
 &\le 2J^{-1}\tau_{\mathrm{src}}(X)+d(|X|-1),\notag\\
 \ell_G(\rho_JX)
 &\le 4J^{-1}\tau_{\mathrm{src}}(X)+2d|X|.             \label{ag-support-projection}
\end{align}
The first line uses \eqref{ag-distance-projection} and
$\ell_{\mathrm{src}}\le2\tau_{\mathrm{src}}$; the second
uses $\ell_G\le2\tau_G$. These conservative constants suffice.

The classical precision has a strict symmetric expansion
$A(V)=A_0+\sum_X B_X(a|_X)$, with rows and columns in $X$.
For some fixed $\sigma_0>0$ and any prescribed finite
$A_{\mathrm{card}}$, the earlier support theorem gives
\begin{equation}\label{ag-source-precision}
 \sup_i\sum_{X\ni i}\sum_{j\in X}
 e^{\sigma_0\tau_{\mathrm{src}}(X)+A_{\mathrm{card}}|X|}
            \|(B_X)_{ij}\|_b\le C_A b.
\end{equation}
The radius may shrink with the reserve. The exponent
$\sigma_0$ can be fixed first because it comes from the
fixed-$L$ linear decay. The original conversion from current
cells to source blocks is already included here.

Use the Gaussian parameters
\begin{equation}\label{ag-Gaussian-exponents}
 \kappa_G=2,\qquad m_G=3,\qquad
 \gamma_G=\kappa_G+2m_G=8.
\end{equation}
The strict Gaussian theorem requires $0<\kappa_G<m_G$;
in particular $m_G=1$ would be invalid when $\kappa_G=2$.
Group precision labels by $X'=\rho_JX$ and their matrix
rows and columns into macrocells. Equation
\eqref{ag-support-projection} shows that its required macro
precision norm at exponent $\gamma_G$ is $O(b)$ if, in $d=4$,
\begin{equation}\label{ag-precision-budget}
 32/J\le\sigma_0,\qquad A_{\mathrm{card}}\ge64.
\end{equation}
The absolute block row sum costs at most $J^d C_A b$:
sum the source row bound over the at most $J^d$ source cells
in one macrocell. Symmetry gives columns. Strict background
support and individual symmetry survive grouping.

This reserve belongs to classical precision coefficients,
where it is available. It is not a new assumption on quantum
activities. The same argument groups the classical remainder
and the local Haar, pivot-Jacobian and gauge multipliers at
their required fixed target weights. Their original classical
or local coefficients admit the necessary finite reserve.
Their small fluctuation bounds remain
$O_{L,J}(gR^3+gR+g^2R^2)$ on the original retained window.

\subsection{Regrouping the actual stationary support family}

On source-block roots and next-current unit background labels,
the matched one-step branch has strict families
\begin{equation}\label{ag-source-stationary}
 a_i(w)=\sum_Xu_{i;X}(w|_X),\qquad
 \sup_i\sum_Xe^{\kappa_0\tau_{\mathrm{src}}(\{i\}\cup X)
                       +A_{\mathrm{st}}|X|}
                  \|u_{i;X}\|_r\le C_{A_{\mathrm{st}}}r,
\end{equation}
for some $\kappa_0>0$ and any prescribed finite
$A_{\mathrm{st}}$, with the corresponding smaller radius.
The output at $i$ is the whole current-background fiber of
one source block.

Group output roots and background labels by $\rho_J$.
For macro root $I$ and macro support $Z$, its $i$-th
source component is $\sum_{X:\rho_JX=Z}u_{i;X}$ for $i\in I$.
Take the maximum over this whole macro output fiber before
the support sum, and bound that maximum by the sum over
$i\in I$. This costs $J^d$. The rooted projection inequality is
\begin{equation}\label{ag-rooted-projection}
 \tau_G(\{\rho_Ji\}\cup\rho_JX)
 \le2J^{-1}\tau_{\mathrm{src}}(\{i\}\cup X)+d|X|.
\end{equation}
Consequently, if $8/J\le\kappa_0$ and $A_{\mathrm{st}}\ge16$,
\begin{equation}\label{ag-macro-stationary}
 \sup_I\sum_Ze^{4\tau_G(\{I\}\cup Z)}\|u_{I;Z}\|_r
 \le J^d C_{A_{\mathrm{st}}}r=:C_{\mathrm{st}}r.
\end{equation}
This is the actual stationary background, with macro spatial
exponent $\kappa_{\mathrm{st}}=4$. Its classical arbitrary
reserve paid \eqref{ag-rooted-projection}. No quantum
reserve was spent. The later Taylor expansion of a general
activity still counts the full dimension $D_{\mathrm{bg}}$.

\subsection{New fluctuation and normalization coefficients}

The macro strict Gaussian theorem now has output reserve
\begin{equation}\label{ag-new-reserve}
 a_G=\min\{1/2,\kappa_G/4\}=1/2.
\end{equation}
Choose $\theta=1/(4D_G)$, so
$\log\mathcal K_\theta\le5/8$. The macro precision norm is
$O(b)$, hence its strengthened $\theta$-dependent smallness
conditions can be imposed by shrinking $b$ with $L,J$ fixed.
The edge constants may be large but remain finite.

The actual cutoff is the product of the retained link/color
balls. A bad macrocell implies $\|\xi_I\|_2>R$, so its joint
tail is bounded by the Euclidean-block tail at the same
threshold $R$. The opposite good-event bound involving the
square root of the number of links cannot improve this bad
threshold. The chart uses the original link supremum $R$.

Two external normalization terms must be included. The local
pivot-delta term
\[
 \log\frac{\mathcal D(V)}{\mathcal D(1)}
       =-\sum_c\operatorname{Tr}\log\pi_c(V)
\]
has fixed source range. After grouping, it has anchored
$O_{L,J}(b)$ norm with any fixed macro support weight,
including $1/2$; a macrocell contains only $O(J^d)$ original
local terms. This follows from its local formula, not from
regrouping an old quantum coefficient family.

For the Gaussian determinant, rerun its chain expansion on
the macro baseline and precision labels. Let $K_\gamma$
and $\nu_\gamma$ be their weighted row norms at $\gamma=8$,
with $q=K_\gamma\nu_\gamma<1$. Then
\begin{equation}\label{ag-det-series}
 \mathcal Q(V)=-\tfrac12\log\det(A(V)A_0^{-1})
   =\tfrac12\sum_{n\ge1}\frac{(-1)^n}{n}
                        \operatorname{Tr}[(C_0B)^n].
\end{equation}
Expand each $B=\sum_YB_Y$, join every entire label $Y$ by
a fixed minimizing tree, and join covariance endpoint pairs
by fixed shortest paths. Their cyclic union is a connected
declared carrier $H$. If $T$ is the sum of these tree/path
lengths and $b_1$ the initial matrix index, then
\begin{equation}\label{ag-det-geometry}
 |H|\le T+1,\qquad z\in H\ \Longrightarrow\ d_G(z,b_1)\le T.
\end{equation}
The graph is connected and has at most $T$ distinct edges;
these inequalities include anchors on added covariance paths.
For $a_Q=1/2$ and $\delta_Q=1$,
\[
 \mathbf1_{\{z\in H\}}e^{a_Q|H|}
 \le e^{1/2}e^{-d_G(z,b_1)}e^{(3/2)T}
 \le e^{1/2}e^{-d_G(z,b_1)}e^{\gamma T}.
\]
The matrix trace is bounded by $D_G$ times the product of
block norms. At fixed $b_1$, alternating the weighted
covariance and precision-label row sums bounds a length-$n$
chain by $D_Gq^n$; dropping the final closing constraint
only enlarges this positive majorant. The precision weight
pays the internal trees since $\tau(Y)\le\ell(Y)$.
Summing $b_1$ gives $S_1=\sum_{z\in\mathbb Z^d}e^{-|z|_1}$.
Thus the connected determinant coefficients satisfy
\begin{equation}\label{ag-det-bound}
 \sup_z\sum_{H\ni z}e^{|H|/2}\|\mathcal Q_H\|
 \le\frac{D_Ge^{1/2}S_1}{2}\sum_{n\ge1}\frac{q^n}{n}
 \le\frac{D_Ge^{1/2}S_1}{2}\frac q{1-q}.
\end{equation}
Each coefficient reads only its precision-label backgrounds,
a subset of $H$. The baseline has no hidden background.
Absolute convergence gives the equality and holomorphy.
This is a fresh macro determinant proof with reserve $1/2$,
not an attempted amplification of an old low-weight family.
Field-independent determinants and powers of $g$ remain
the separately tracked vacuum normalization.

Together, the new Gaussian and normalization coefficients
have reserve $1/2$ and a norm $\epsilon_{\mathrm{new}}$
bounded by $O_{L,J}(b)$ plus the stated fluctuation-log bound.
This concerns their support weight, not an all-scale
smallness or contraction assertion.

\subsection{Inherited interactions still project by the ratio L}

Assume inherited connected activities on $\Lambda_Q$ satisfy
\begin{equation}\label{ag-inherited-input}
 \sup_i\sum_{H\ni i}e^{a_0|H|}\|E_H\|_b
       \le\epsilon_{\mathrm{old}},\qquad a_0=1/16.
\end{equation}
The source fluctuation map still has ratio $L$. On the
independent product background $\|a\|<b_-<b$, its exact
increment obeys
\begin{equation}\label{ag-inherited-increment}
 \|\Delta E_H\|\le u\|E_H\|_b,\qquad
 u=\frac{C_{\mathrm{ch}}(L)|g|R}
        {b-b_- - C_{\mathrm{ch}}(L)|g|R},
\end{equation}
when the chart and positive denominator conditions hold.
Directional Cauchy gives this estimate without a
$J$-dependent coordinate count.

Nonpivot outputs are bond-local. A pivot output reads two
source $L$-blocks, which either lie in one macrocell or in
adjacent macrocells. As both the source and quantum tilings
refine the common $JL$ macrocell, its declared carrier obeys
\begin{align}
 Y(H)&\text{ is connected},&
 |Y(H)|&\le\chi|\pi_LH|,\quad\chi=d+1,\notag\\
 \mathbf1_{\{B\in Y(H)\}}
 &\le\sum_{i\in H}\mathbf1_{\{d_G(B,\pi_Li)\le1\}}.
                                                               \label{ag-pivot-geometry}
\end{align}
The safe anchor factor is $N_{\mathrm{near}}L^d$, with
$N_{\mathrm{near}}=2d+1$; there is no extra $J^d$.
Connected-carrier compression in the quantum graph gives
\begin{equation}\label{ag-quantum-compression}
 |\pi_LH|\le2^d+(2^d/k_L)|H|,\qquad
 k_L=1+\lfloor(L-3)/4\rfloor.
\end{equation}
The target Gaussian multiplier exponent is at most
\begin{equation}\label{ag-required-weight}
 B_{\mathrm{req}}=4+\log\mathcal K_\theta+\kappa_G
                  \le53/8.
\end{equation}
Consequently the inherited interaction weight is paid if
\begin{equation}\label{ag-L-choice}
 \frac{\chi2^d}{k_L}\frac{53}{8}\le\frac1{16}.
\end{equation}
In dimension four this asks for $k_L\ge8480$. For example
$L=33921$ gives equality and $L=33925$ a strict margin.
These crude numbers demonstrate a finite choice, not a
useful numerical coupling range.

After grouping equal carriers, let
$Q_Y=\sum_{H:Y(H)=Y}\Delta E_H$ and
$V_{\max}=uN_{\mathrm{near}}L^d\epsilon_{\mathrm{old}}$.
The exact multiplier estimate yields
\begin{equation}\label{ag-inherited-multiplier}
 \epsilon_{\mathrm{int}}
 \le e^{V_{\max}}u e^{(53/8)\chi2^d}
             N_{\mathrm{near}}L^d\epsilon_{\mathrm{old}}.
\end{equation}
Indeed $\|Q_Y\|\le V_{\max}$ and
$\|e^{Q_Y}-1\|\le e^{\|Q_Y\|}\|Q_Y\|$;
\eqref{ag-pivot-geometry} and
\eqref{ag-quantum-compression} then give the anchored sum.
The base and inherited factors may share a support as
distinct labels. Under the previously proved quarter
budgets their connected-log difference is at most
$4C_{\mathrm{out}}\epsilon_{\mathrm{int}}$.
The old fluctuation-independent sum $\sum_HE_H(a)$ stays
outside this integral.

\subsection{Transport of both kinds of background term}

Choose a positive common radius $b_{\mathrm{com}}$ no
larger than $b_-$ and every domain of the new Gaussian,
determinant and delta coefficients. In particular
$b_{\mathrm{com}}\le b$; restricting the old activities to
this smaller ball preserves their input norm. The inherited
outer radius $b$ cannot silently be assigned to new
coefficients constructed only on $b_-$.

Use \eqref{ag-macro-stationary} and the transport parameters
\begin{equation}\label{ag-transport-budget}
 \beta=1/8,\quad a_{\mathrm{out}}=\mu=1/16,\quad
 \delta_{\mathrm{tr}}=1,\qquad
 t=\frac{C_{\mathrm{st}}r}{b_{\mathrm{com}}}
                      \le\frac1{16D_{\mathrm{bg}}}.
\end{equation}
Then $\beta+\delta_{\mathrm{tr}}<\kappa_{\mathrm{st}}=4$
and $D_{\mathrm{bg}}\log(1/(1-t))\le1/8$.
New Gaussian and normalization coefficients already have
carriers in $\Lambda_G$ and reserve $1/2$. Identity-projection
transport returns reserve $1/16$ and norm at most
\begin{equation}\label{ag-new-transport}
 S_1\epsilon_{\mathrm{new}}.
\end{equation}
Their output carriers are $J$-blocks of the next current
lattice, exactly the type assumed at the next stage.

The old background term has carriers in $\Lambda_Q$ and
must instead use $\pi_L$. Each of its $J$-current output
fibers is a restriction of the containing $JL$ macro fiber,
so it inherits the rooted bound anchored at $\pi_L(i)$.
Compressed transport uses $c_0=2^d,c_1=2^d/k_L$ and
the coordinate count $D_{\mathrm{in}}$. For $d=4$,
\begin{equation}\label{ag-old-transport-budget}
 c_1\beta+D_{\mathrm{in}}\log\frac1{1-t}
 \le\frac2{k_L}+\frac1{8L^d}<\frac1{16},
\end{equation}
because $D_{\mathrm{in}}/D_{\mathrm{bg}}=L^{-d}$.
It too returns reserve $1/16$, with norm at most
\begin{equation}\label{ag-old-transport}
 e^{\beta2^d}L^d S_1\epsilon_{\mathrm{old}}
       =e^2L^d S_1\epsilon_{\mathrm{old}}.
\end{equation}
Its field-independent constants are tracked separately as
vacuum terms. No reserve beyond $1/16$ was assumed for it.
Treating it as if it already had the Gaussian output reserve
$1/2$ would have been incorrect.

Thus the full update has the same connected-carrier exponent
$a_0=1/16$ and a one-step bound
\begin{equation}\label{ag-weight-recurrence}
 \epsilon_{\mathrm{next}}
 \le e^2L^dS_1\epsilon_{\mathrm{old}}
                     +S_1\epsilon_{\mathrm{new}},
\end{equation}
on the common smaller field domain. The inherited
fluctuation effect is already included in
$\epsilon_{\mathrm{new}}$ through
\eqref{ag-inherited-multiplier}. This is recovery of a
function-space weight, not contraction of its norm.

\subsection{A finite choice of J and the remaining induction boundary}

After fixing $L$ in \eqref{ag-L-choice}, choose $J=L^j$
large enough that
\begin{equation}\label{ag-J-choice}
 8/J\le\gamma_0,\qquad32/J\le\sigma_0,\qquad
 8/J\le\kappa_0.
\end{equation}
Every exponent on the right is positive for that fixed $L$,
uniformly in $M$. Choose classical reserves at least $64$
for precision and $16$ for the stationary family, and any
additional finite amount needed for classical interaction
grouping. Their established constructions permit this at
a smaller radius. Hence the choices have a finite solution
without any large-$L$ asymptotic estimate. The quantum
family always starts on $J$-current cells; no extra quantum
reserve is spent to amplify a macro spatial exponent.

The analytic radius is still a real obstruction. In particular,
\begin{equation}\label{ag-radius-loss}
 r\le\frac{b_{\mathrm{com}}}{16D_{\mathrm{bg}}C_{\mathrm{st}}},
 \qquad C_{\mathrm{st}}=J^dC_{A_{\mathrm{st}}},
\end{equation}
must be imposed along with all chart and precision domains.
An incoming radius may itself have been produced by a
previous step. This theorem does not turn
\eqref{ag-radius-loss} into a self-map, nor compare it with
the required running-coupling domains. The smoother pre-gauge
representative and its analytic-germ extension have separate
integral and activity-norm interfaces.

Similarly, \eqref{ag-weight-recurrence} has a coefficient
larger than one, and the normalization terms are $O(b)$,
not automatically $O(g)$. Uniform amplitude control requires
vacuum, relevant and marginal extraction and a coupling-flow
estimate. Large-field pieces, mixed histories and observable
bounds are unchanged obligations.

The result is therefore specific: fixed $L,J$ and the exact
source coordinates preserve a fixed positive connected
support exponent through inherited insertion, normalized
Gaussian integration, external normalization and stationary
transport. This resolves the earlier circular support-weight
test while leaving the analytic-domain and amplitude induction
explicitly open.
